\documentclass[10pt,a4paper]{amsart}
\usepackage[margin=19mm]{geometry}
\usepackage{amsmath,amssymb,mathtools}
\usepackage[scr=rsfs,cal=euler]{mathalfa}
\usepackage{xfrac}
\usepackage[unicode,hidelinks,bookmarks=false]{hyperref}
\newcommand{\ie}{i.e.\ }
\newcommand{\cf}{cf.\ }
\newcommand{\resp}{respectively\ }
\newcommand{\Subsection}{\subsection}
\providecommand{\nobreakdash}{\nobreak\hbox{-}}
\setlength{\emergencystretch}{2em}
\multlinegap0pt
\usepackage{bm}
\usepackage{bbm}
\usepackage{dsfont}
\usepackage{xspace}
\usepackage{cancel}
\usepackage{mathtools}
\usepackage{amsthm}
\makeatletter
\@ifundefined{theorem}{\newtheorem{theorem}{Theorem}}{}
\@ifundefined{lemma}{\newtheorem{lemma}[theorem]{\bf Lemma}}{}
\makeatother
\newcommand\cP{{\mathcal P}}
\title[Avoiding configurations of small size in the square grid]{Avoiding configurations of small size in the square grid}
\author{M\'at\'e J\'anosik, Art\'ur N\'ador, Zolt\'an L\'or\'ant Nagy, L\'aszl\'o Bence Simon}
\date{Source: arXiv:2601.14465v2 (CC BY 4.0)}
\begin{document}
\maketitle

\noindent\emph{Source and licence.} Avoiding configurations of small size in the square grid. Version arXiv:2601.14465v2 (CC BY 4.0), distributed under the Creative Commons Attribution 4.0 International licence, as recorded in the arXiv licence metadata for this exact version and shown on the abstract page https://arxiv.org/abs/2601.14465v2. Full rights notice: licenses/arxiv-2601.14465-CC-BY-4.0.txt.\par\smallskip

\noindent\textit{Editorial scope.} Counted unit: the Lemma whose source label is \texttt{isocount} in the article (source0.tex lines 194--196, source label \texttt{isocount}), delivered together with the article's own complete original proof of it (source0.tex lines 197--243, one \texttt{proof} environment of 47 lines). No mathematics is written, repaired or re-derived: statement and proof are the article's own text line for line. Required context retained uncounted: none. Every definition, display and label of those lines is retained unchanged. Omitted from this package and not redistributed: 1--193, 244--680. Nothing inside a retained excerpt is omitted or altered: every retained body is delivered exactly as the article's lines are written, so no omission receipt of any kind accompanies this package. Citations: the retained excerpts carry no citation command at all, so no bibliography entry of the article is transcribed and no reference is redistributed. Reference adaptation: none was needed and none was made; every reference inside the delivered unit resolves inside it, so no unresolved reference or citation is produced. Printed slips kept literal: no printed word, symbol, hypothesis or number of the counted statement or of its proof was altered. Layout and class: the article's own class is \texttt{article} with packages including graphicx, inputenc, amsmath, mathtools, amssymb, amsthm, amscd, epic, eepic, url, color, makebox, comment, enumerate; the wrapper is the lane's pinned \texttt{amsart} editorial wrapper at 10pt on A4, which restates the theorem-like environments the retained text needs and adds the article's own macro definitions that the retained text needs (\texttt{\textbackslash cP}), with extra wrapper packages (bm, bbm, dsfont, xspace, cancel, mathtools). The delivered numbering is the unit's own. Assets: no retained span contains a figure environment, a graphics-inclusion command, a PDF-page inclusion, a table or a TikZ picture; no image file is copied into this package, the package declares no asset dependency and no image file is redistributed with it. Licence: version arXiv:2601.14465v2 of this article is distributed under the Creative Commons Attribution 4.0 International licence, as recorded in the arXiv licence metadata for this exact version and shown on the abstract page https://arxiv.org/abs/2601.14465v2. A full rights notice is delivered at licenses/arxiv-2601.14465-CC-BY-4.0.txt. Characters: every retained excerpt is pure ASCII, so the delivered engine, XeLaTeX with its default Latin Modern text font, reports no missing character.
	\begin{lemma}\label{isocount}
	The number $T(\cP)$ of isosceles triangles in the point set $\cP:=S\times S$ satisfies $T(\cP)=O(|S|^{\beta}).$
	\end{lemma}

	\begin{proof}
	Let $S\subset\{1,\dots,n\}$ be a Behrend-type construction, $m:=|S|$, and set $\cP:=S\times S\subset\mathbb Z^2$.  We will count ordered triples $(A,B,C)\in \binom {\cP}{3}$, where $B\ne C$ form the base of an isosceles triangle, and $A$ is the peak, i.e.,  equidistant from $B$ and $C$. Counting ordered triples only changes $T$ by a factor of $2$ as no equilateral triangle lie in $\mathbb Z^2$.
	
	The isosceles condition can be  written as
	\[
	(x_B-x_A)^2+(y_B-y_A)^2=(x_C-x_A)^2+(y_C-y_A)^2
	\]
	in terms of the coordinates of $A,B,C$ which is equivalent to 
	\[
	2(x_B-x_C)x_A + 2(y_B-y_C)y_A = x_B^2-x_C^2 + y_B^2-y_C^2.
	\]
	Let us define
	\[
	a:=x_B-x_C,\qquad b:=y_B-y_C,\qquad c:=\frac{x_B^2-x_C^2+y_B^2-y_C^2}{2}.
	\]
	With this notion, the coordinates of $A$ must satisfy the linear equation of the bisector
	\begin{equation}\label{line}
	a x_A + b y_A = c,
	\end{equation}
	and the peak must be a grid  point on the bisector lying in $\cP=S\times S$.
	Let us 	fix a nonzero integer pair $(a,b)\neq(0,0)$ and an integer $c$. The integer solutions to
	\[
	a x + b y = c
	\]
	exist only if $\gcd(a,b)$ divides $c$, and in that case the solutions can be written as 
	
	\[
	x = x_0 + \frac{b}{\gcd(a,b)}t,\qquad y = y_0 - \frac{a}{\gcd(a,b)} t,\qquad t\in\mathbb Z,
	\] where  $(x_0,y_0)$ is an integer solution of \eqref{line}. Note that while counting isosceles triangles on base $BC$ in $S\times S$, we can also assume that neither $a$, nor $b$ is equal to $0$, as $S$ is $3$-AP-free.
	
	
	Thus the possible peak points w.r.t. the pair $B, C$ are parametrized by  parameter $t$, and forming an arithmetic progression in the $x$-coordinate with step $|b|/\gcd(a,b)$ (correspondingly in the $y$-coordinate with step $|a|/\gcd(a,b)$).
	
	Hence, if we set
	\[
	m:=\frac{|b|}{\gcd(a,b)},
	\]
	then the number of peaks from $\cP=S\times S$ on this line is at most the number of integers $t$ for which both coordinates satisfy
	\[
	x_0 + \frac{b}{\gcd(a,b)}t\in S,\qquad y_0 - \frac{a}{\gcd(a,b)}t\in S.
	\]
	In particular, it is bounded by the maximum size $M(m)$ of an intersection of $S$ with a single residue class modulo $m$. Let us denote by $N(m)$ the number of pairs $(B,C)$ whose bisector parameters $a,b$ yield modulus $m$. Then in order to count the number of isosceles triangles, it is enough to prove an upper bound on $\sum_m N(m)M(m)$ as we clearly have $\sum_m N(m)M(m)\ge T(\cP)$.
	
	First we show the upper bound $N(m)\le 2n^2|S|^2/m$. Recall that $m=\frac{|b|}{\gcd(a,b)}$ where $a=x_B-x_C$ and $ b=y_B-y_C $, thus the number of $4$-tuples $(x_B, x_C, y_B, y_C)$ yielding modulus $m$ is at most $|S|^2$ times the number of solutions for $m=\frac{b}{\gcd(a,b)}$ among the pairs of integers $(a,b)\in [-(n-1), (n-1)]^2$. Since $|b|$ must be a positive multiple of $m$ which is smaller than $n$, the number of such solutions is bounded from above by $\frac{n}{m}\cdot 2n$, which in turn proves the claim.\\
	Our next aim is to show that $M(m)\leq (1+o(1))\frac{|S|}{m}$ holds as $n\to \infty,$ i.e., that $S$ is almost equidistributed among the residue classes $\pmod m.$
	
	\end{proof}

\end{document}
